\documentclass[12pt,letterpaper]{article}

\usepackage[letterpaper,left=1.18in,right=1.18in,top=1.0in,bottom=0.9in,footskip=0.35in]{geometry}
\usepackage{amsmath,amssymb,mathtools,bm}
\usepackage{graphicx}
\usepackage{booktabs}
\usepackage{caption}
\usepackage{enumitem}
\usepackage{titlesec}
\usepackage{xcolor}
\usepackage[colorlinks=true,linkcolor=blue,citecolor=blue,urlcolor=blue,linktocpage=true]{hyperref}

\setlength{\parindent}{1.2em}
\setlength{\parskip}{0pt}
\linespread{1.05}

\numberwithin{equation}{section}

\captionsetup{
    labelfont=normalfont,
    textfont=normalfont,
    labelsep=period,
    justification=centering
}

\titleformat{\section}
    {\normalfont\Large\bfseries}
    {\thesection}
    {1.15em}
    {}
\titleformat{\subsection}
    {\normalfont\large\bfseries}
    {\thesubsection}
    {1.0em}
    {}
\titleformat{\subsubsection}
    {\normalfont\normalsize\bfseries}
    {\thesubsubsection}
    {0.9em}
    {}
\titlespacing*{\section}{0pt}{2.6ex plus 0.5ex minus 0.2ex}{1.8ex plus 0.2ex}
\titlespacing*{\subsection}{0pt}{2.2ex plus 0.4ex minus 0.2ex}{1.0ex plus 0.2ex}
\titlespacing*{\subsubsection}{0pt}{1.8ex plus 0.3ex minus 0.2ex}{0.8ex plus 0.1ex}

\makeatletter
\def\ps@plain{%
    \let\@mkboth\@gobbletwo
    \let\@oddhead\@empty
    \let\@evenhead\@empty
    \def\@oddfoot{\hfil-- \thepage\ --\hfil}%
    \let\@evenfoot\@oddfoot
}
\pagestyle{plain}

\renewcommand{\tableofcontents}{%
    \section*{\contentsname}%
    \@starttoc{toc}%
}
\renewcommand{\l@section}[2]{%
    \ifnum \c@tocdepth >\z@
        \addpenalty\@secpenalty
        \vskip 0.8em \@plus\p@
        \setlength\@tempdima{1.65em}%
        \begingroup
            \parindent \z@
            \rightskip \@pnumwidth
            \parfillskip -\@pnumwidth
            \leavevmode
            \bfseries
            \advance\leftskip\@tempdima
            \hskip -\leftskip
            #1\nobreak\hfil \nobreak\hb@xt@\@pnumwidth{\hss #2}\par
        \endgroup
    \fi}
\renewcommand{\l@subsection}{\@dottedtocline{2}{1.65em}{2.7em}}
\renewcommand{\l@subsubsection}{\@dottedtocline{3}{4.35em}{3.4em}}
\renewcommand{\@dotsep}{10000}
\makeatother

\newcommand{\notestitle}{Testing Pattern-Mediated\\Indirect Evolutionary Rescue}
\newcommand{\notessubtitle}{A predator--prey reaction--diffusion research report}
\newcommand{\notesauthor}{Quoc Trong Do}

\newenvironment{lecturequote}[1]
{\def\lecturequoteauthor{#1}\begin{quote}\itshape}
{\par\medskip\hfill \lecturequoteauthor\end{quote}}

\begin{document}

\hypersetup{pageanchor=false}

\pagenumbering{roman}
\begin{center}
{\sffamily\bfseries\LARGE \notestitle\par}
\vspace{0.18in}
{\sffamily\bfseries\large \notessubtitle\par}
\vspace{0.15in}
{\bfseries \notesauthor\par}
\end{center}
\vspace{0.15in}
\noindent\rule{\linewidth}{0.5pt}
\vspace{0.2in}

\hypersetup{pageanchor=true}
\tableofcontents
\clearpage

\pagenumbering{arabic}
\setcounter{page}{1}

\section{Introduction}

\begin{lecturequote}{Philip Anderson}
This process of model building is the art of discarding all but the essentials and
focusing on a model simple enough to do the job, but not too hard to see all the way
through.
\end{lecturequote}

\subsection{Basic Biological Concepts}

We begin by fixing the biological language. A \emph{population} is a group of
individuals of the same type or species. In mathematical models, its abundance may be
tracked as a function of time, such as $X(t)$, or as a density in space and time, such
as $X(x,t)$.

A \emph{prey} population is a population that can be consumed by another population.
A \emph{predator} population feeds on prey, and the interaction in which predators
consume prey is called \emph{predation}. Predation has opposite immediate effects on
the two populations: it decreases prey abundance but can increase predator abundance by
providing energy for survival and reproduction.

A \emph{trait} is a biological property of individuals that affects growth, survival,
movement, or interactions with other species. In this report, the key prey trait is
defence against predation. We distinguish between two prey types:
\begin{equation}
    U=\text{undefended prey},
    \qquad
    D=\text{defended prey}.
\end{equation}
Undefended prey are easier for predators to consume, but they are assumed to grow or
compete better. Defended prey are harder to consume, but defence is costly. This is the
defence--growth trade-off that motivates the extended model.

An \emph{eco-evolutionary} model is one in which ecological variables, such as
population density, interact with evolutionary variables, such as trait composition.
In our setting, the predator may be rescued not because it evolves, but because the prey
community shifts toward a more edible composition.

Finally, a \emph{spatial} model tracks where populations are located. This matters
because predator persistence depends not only on how much edible prey exists, but also
on whether predators and edible prey overlap in space. A \emph{Turing pattern} is a
spatial pattern generated by reaction--diffusion dynamics when diffusion destabilizes
an otherwise stable homogeneous state \cite{turing1952}. The central question of this
report is whether such spatial patterns can change the threshold for predator rescue.

\subsection{Literature Positioning and Research Gap}

The starting point is the broader problem of \emph{evolutionary rescue}: a population
under environmental stress may avoid extinction if evolutionary change raises its mean
fitness fast enough \cite{gomulkiewicz1995}. Most evolutionary-rescue theory focuses
on the threatened population itself. In multispecies communities, however, a focal
species can also be rescued indirectly by evolutionary change in an interacting species.
Yamamichi and Miner introduced this mechanism as \emph{indirect evolutionary rescue}:
prey evolution can prevent predator extinction when the prey community shifts toward a
more edible composition \cite{yamamichi2015}. Hermann and Becks later tested this idea
experimentally in a rotifer--alga system, showing that changes in prey genotype
frequency can rescue a predator from extinction \cite{hermann2022}.

This report builds on that predator--prey setting but focuses on a spatial question.
Defended and undefended prey are a natural way to represent the key trait axis because
defence typically reduces predation risk but carries a growth or competitive cost.
Experiments with rotifer--algal systems show that the form of this trade-off can
strongly alter eco-evolutionary dynamics \cite{kasada2014}, and recent work shows that
the range of prey trait variation affects the evolutionary contribution to predator
growth rates \cite{hermann2024}. These studies motivate treating prey composition as a
dynamic variable rather than as a fixed ecological background.

The closest spatial precedent is the reaction--diffusion model of Roy et al., who
studied defended and undefended prey, eco-evolutionary feedback, coexistence,
persistence, extinction, and Turing pattern formation in a spatial predator--prey system
\cite{roy2026}. Their work shows that spatio-temporal patterns can arise from the
interaction between ecological and evolutionary components. The present report adds a
more targeted threshold question: instead of asking only when patterns occur or which
states persist, we ask whether spatial pattern formation changes the maximum predator
mortality compatible with indirect rescue. The quantity
$\Delta m_c=m_c^{\mathrm{PDE}}-m_c^{\mathrm{ODE}}$ is therefore the operational
comparison that distinguishes this report from the broader Roy et al. pattern analysis.
Thus ``pattern-mediated indirect evolutionary rescue'' is used here as a proposed
mechanism to test, not as an established subfield.

The classical Lotka--Volterra model provides the baseline language for this question.
It is not realistic enough for the final problem, but it clearly displays the
predator--prey feedback that all later models retain: prey support predators, while
predators suppress prey. We therefore begin with that baseline, identify its
limitations, and then extend it to include prey heterogeneity, predator stress, spatial
movement, and a rescue-threshold comparison.

\subsection{Core Mathematical Problem}

The study is not merely about finding Turing patterns. Turing instability is a
mechanism that can create spatial structure, but the central mathematical problem is to
determine whether such structure changes the predator mortality threshold for indirect
evolutionary rescue.

The mathematical object of interest is the difference between the predator persistence
threshold in the well-mixed system and the corresponding threshold in the spatial
reaction--diffusion system. Let $m_c^{\mathrm{ODE}}$ denote the largest predator
mortality for which the non-spatial model maintains predator persistence, and let
$m_c^{\mathrm{PDE}}$ denote the corresponding threshold for the spatial model. The
central quantity is
\begin{equation}
    \Delta m_c=m_c^{\mathrm{PDE}}-m_c^{\mathrm{ODE}}.
\end{equation}
If $\Delta m_c>0$, spatial patterning promotes indirect evolutionary rescue. If
$\Delta m_c<0$, spatial patterning inhibits rescue. If $\Delta m_c$ changes sign
across parameter regimes, then the effect of pattern formation is conditional.

\begin{center}
\emph{Do Turing-type spatial patterns in defended and undefended prey expand or shrink
the mortality range over which predators persist?}
\end{center}

\section{Predator--Prey Dynamics}

\subsection{The Basic Ecological Problem}

A predator--prey system is the simplest setting in which the growth rate of one
population depends directly on the abundance of another. Let $X(t)$ denote prey
density and $Y(t)$ predator density. If prey are abundant, predators have more
opportunities to feed and reproduce. If predators are abundant, prey suffer more
mortality. This creates the feedback loop
\[
    \text{more prey} \longrightarrow \text{more predator growth}
    \longrightarrow \text{more prey mortality}.
\]
After prey decline, predators may also decline because their food supply has been
reduced. Predator decline then releases prey from predation pressure. The classical
predator--prey problem is to turn this verbal feedback into a dynamical system.

\subsection{Historical and Empirical Motivation}

The equations are named after Alfred J. Lotka and Vito Volterra, who arrived at the
same mathematical structure from different directions. Lotka introduced related
equations in the context of chemical reactions and later used them in population
dynamics. Volterra developed the predator--prey version after discussions with the
marine biologist Umberto D'Ancona, who had studied fish catches in the Adriatic Sea.
During the First World War, fishing effort was greatly reduced, yet the proportion of
predatory fish in the catch increased. Volterra's model gave a simple mathematical
mechanism for why reduced fishing could increase the average abundance of predators
relative to prey \cite{murray2002,kot2001}.

The model is also associated with the famous hare--lynx records of the Hudson's Bay
Company. The recorded numbers of snowshoe hare and lynx pelts show repeated
oscillations that resemble predator--prey cycles. These data are not a clean laboratory
experiment: pelt records are affected by hunting effort, climate, food availability,
disease, and many other variables. Still, they remain a useful empirical picture of the
kind of delayed predator--prey oscillation that the model tries to capture. Cleaner
cyclic behaviour can be observed in controlled systems, such as laboratory
rotifer--algae predator--prey experiments, where environmental conditions and species
composition are more tightly constrained \cite{kot2001,edelstein2005}.

These examples should not be taken as proof that the Lotka--Volterra equations are a
complete model of real ecosystems. Rather, they explain why the model became central:
it is simple enough to analyze, but rich enough to produce a recognizable ecological
story.

\subsection{Classical Lotka--Volterra Baseline}

The Lotka--Volterra model is the standard first mathematical answer to the
predator--prey problem \cite{murray2002,kot2001,edelstein2005}. It is built from
four assumptions:

\begin{enumerate}[leftmargin=2.2em]
    \item The system is well mixed, so predator--prey encounters are proportional to
    the product $XY$.
    \item In the absence of predators, prey grow exponentially.
    \item In the absence of prey, predators decline exponentially.
    \item The environment is constant, and there is no trait change, adaptation, age
    structure, space, or stochastic extinction.
\end{enumerate}

These assumptions are strong, but they give the cleanest possible expression of
predator--prey feedback. Combining exponential prey growth, predator mortality, and
mass-action predation gives
\begin{equation}
\boxed{
\begin{aligned}
    \frac{dX}{dt} &= \alpha X-\beta XY,\\
    \frac{dY}{dt} &= \delta XY-\gamma Y.
\end{aligned}}
\label{eq:lotka-volterra}
\end{equation}
Here $\alpha$ is the prey growth rate, $\beta$ is the attack coefficient on prey,
$\delta$ converts prey consumption into predator growth, and $\gamma$ is predator
mortality.

The per capita form is
\begin{equation}
    \frac{1}{X}\frac{dX}{dt}=\alpha-\beta Y,
    \qquad
    \frac{1}{Y}\frac{dY}{dt}=\delta X-\gamma.
\end{equation}
This form shows two thresholds. Prey increase when predator density is below
$\alpha/\beta$, and predators increase when prey density is above $\gamma/\delta$.
Equivalently, the non-trivial nullclines are
\begin{equation}
    \dot X=0 \Longleftrightarrow Y=\frac{\alpha}{\beta},
    \qquad
    \dot Y=0 \Longleftrightarrow X=\frac{\gamma}{\delta}.
\end{equation}
They intersect at the coexistence equilibrium
\begin{equation}
    (X^*,Y^*)=\left(\frac{\gamma}{\delta},\frac{\alpha}{\beta}\right),
\end{equation}
besides the extinction equilibrium $(0,0)$.

The coexistence equilibrium is a center in the classical model, leading to neutral
closed cycles in the positive quadrant. This cyclic behaviour is the famous qualitative
prediction: prey increase, predators respond, prey are suppressed, predators decline,
and prey recover. The result is mathematically elegant, but fragile. Once the model is
perturbed, for example by adding density-dependent prey growth, neutral cycling
generally disappears \cite{murray2002,kot2001}.

\subsection{Why the Classical Model Is Not Enough}

The Lotka--Volterra model is valuable because it is transparent. Every term has a
clear biological meaning, and the predator--prey feedback is easy to identify.
However, it is too simple for the problem we ultimately want to study.

First, prey grow exponentially when predators are absent. This means the model has no
resource limitation and no carrying capacity. A natural correction is logistic prey
growth, where reproduction slows as prey density approaches an ecological ceiling.
This correction is not a cosmetic detail: in the standard logistic predator--prey
extension, the coexistence state becomes an attracting equilibrium rather than a
neutrally stable center.

Second, the closed orbits imply that populations can recover from arbitrarily small
positive densities. This is a deterministic feature of the equations, not a robust
ecological fact. In real populations, very low abundance can mean demographic
stochasticity, Allee effects, genetic loss, or local extinction.

Third, all prey are identical. There is no defended prey type and no undefended prey
type. Therefore the model cannot represent a defense--growth trade-off, prey trait
composition, or evolutionary sorting.

Fourth, the system is well mixed. The variables $X(t)$ and $Y(t)$ describe total
population densities, not spatial distributions. The model cannot ask whether
predators and edible prey overlap in space, or whether diffusion-driven patterns help
or harm predator persistence.

Fifth, there is no explicit rescue threshold. The predator either persists in the
closed-orbit dynamics or is absent from the model. To study indirect evolutionary
rescue, we need a parameter such as predator mortality and a criterion for whether the
predator population avoids extinction.

These limitations are not failures of the Lotka--Volterra model. They tell us what
must be added. The rest of this report keeps the same modelling discipline, but adds
the ingredients required for the research question: prey heterogeneity, predator
stress, spatial movement, and a comparison between non-spatial and spatial persistence
thresholds.

\section{The Extended Lotka--Volterra Model}

\subsection{From One Prey Type to Two}

The Lotka--Volterra system has one prey population and one predator population. Our
problem requires one additional biological layer: prey are not all equivalent. We split
the prey into two types:
\begin{equation}
    U(t)=\text{undefended prey}, \qquad
    D(t)=\text{defended prey},
\end{equation}
and keep
\begin{equation}
    P(t)=\text{predator}.
\end{equation}
The total prey density is
\begin{equation}
    N(t)=U(t)+D(t).
\end{equation}

The distinction between $U$ and $D$ is defined by a trade-off. Undefended prey are
more edible but grow faster. Defended prey are less edible but pay a growth or
competitive cost. We write this as
\begin{equation}
    r_U>r_D,
    \qquad
    a_U>a_D.
\end{equation}
Here $r_U$ and $r_D$ are prey growth rates, while $a_U$ and $a_D$ are predator attack
rates on undefended and defended prey respectively. This two-type reduction is a
minimal representation of the defence--growth trade-off observed in rotifer--algal
systems and related eco-evolutionary experiments \cite{kasada2014,hermann2024}.

The model will be used in two stages. First, the ODE model defines the well-mixed
baseline and the threshold $m_c^{\mathrm{ODE}}$. Second, the PDE model adds diffusion
and defines $m_c^{\mathrm{PDE}}$. The spatial model is therefore evaluated relative to
the non-spatial baseline, not as an isolated pattern-forming system.

\subsection{A Non-Spatial Baseline}

The first extension of Lotka--Volterra is to replace unlimited prey growth by logistic
growth and to let the predator consume two prey types. A minimal non-spatial model is
\begin{equation}
\boxed{
\begin{aligned}
    \frac{dU}{dt}
    &= r_U U\left(1-\frac{U+D}{K}\right)-a_U U P,\\
    \frac{dD}{dt}
    &= r_D D\left(1-\frac{U+D}{K}\right)-a_D D P,\\
    \frac{dP}{dt}
    &= e(a_U U+a_D D)P-mP.
\end{aligned}}
\label{eq:two-prey-sorting}
\end{equation}
This is the direct analogue of Lotka--Volterra for our setting. The prey equations
contain logistic growth and predation loss. The predator equation says that predators
increase by converting consumed prey into predator biomass or offspring, and decrease
at mortality rate $m$.

The parameter $m$ is important because it will become the environmental stress
parameter. Increasing $m$ makes predator persistence harder. In a rescue problem, the
central question is not merely whether the equations oscillate or equilibrate, but
whether the predator avoids extinction when $m$ is high.

\subsection{Evolutionary Sorting}

Although \eqref{eq:two-prey-sorting} contains no explicit mutation term, it still
contains an evolutionary component if $U$ and $D$ represent heritable prey types
already present in the population. The evolutionary variable is the prey composition:
\begin{equation}
    q_U=\frac{U}{U+D},
    \qquad
    q_D=\frac{D}{U+D}.
\end{equation}
In this minimal model, evolution is represented at the level of prey-type frequency
dynamics rather than allele-level population genetics.
To see how selection acts, compare the per capita growth rates of the two prey types.
For positive $U$ and $D$,
\begin{align}
    \frac{d}{dt}\log\left(\frac{U}{D}\right)
    &=
    \frac{1}{U}\frac{dU}{dt}
    -
    \frac{1}{D}\frac{dD}{dt} \notag\\
    &=
    (r_U-r_D)\left(1-\frac{U+D}{K}\right)
    -
    (a_U-a_D)P.
\end{align}
This equation is the simplest expression of the rescue mechanism. When predator
density is high, the second term is large and defended prey are favoured. When predator
density falls, selection against undefended prey weakens, and the faster growth of
undefended prey can increase the edible component of the prey community. Predator
recovery can then occur indirectly through prey compositional change
\cite{yamamichi2015}.

\subsection{A Coexistence Caveat}

There is a mathematical point that matters for the rest of the project. In the pure
sorting model \eqref{eq:two-prey-sorting}, a positive equilibrium with
$U^*>0$, $D^*>0$, and $P^*>0$ is not generic. If all three populations are positive at
equilibrium, then the first two equations imply
\begin{equation}
    r_U\left(1-\frac{U^*+D^*}{K}\right)=a_U P^*,
    \qquad
    r_D\left(1-\frac{U^*+D^*}{K}\right)=a_D P^*.
\end{equation}
Therefore a coexistence equilibrium with both prey types requires
\begin{equation}
    \frac{r_U}{a_U}=\frac{r_D}{a_D}.
\end{equation}
This is a fine-tuning condition, not a robust biological assumption. Thus
\eqref{eq:two-prey-sorting} is useful for studying transient sorting and rescue, but it
is not the safest baseline for a Turing analysis around a homogeneous coexistence
state.

For this reason, the working model includes a small conversion term between prey
types. This is a phenomenological modelling device: it keeps both prey states present
so that a homogeneous coexistence equilibrium can be used for local stability and
Turing analysis. Biologically, such terms could be motivated by mutation, reversible
phenotypic switching, inducible defence, or unresolved standing variation, but those
mechanisms are not equivalent and are not distinguished in this minimal model. Let
$\mu_{UD}$ be the rate from undefended to defended prey, and $\mu_{DU}$ the rate from
defended to undefended prey. The non-spatial working model is
\begin{equation}
\boxed{
\begin{aligned}
    \frac{dU}{dt}
    &= r_U U\left(1-\frac{U+D}{K}\right)-a_U U P
    +\mu_{DU}D-\mu_{UD}U,\\
    \frac{dD}{dt}
    &= r_D D\left(1-\frac{U+D}{K}\right)-a_D D P
    +\mu_{UD}U-\mu_{DU}D,\\
    \frac{dP}{dt}
    &= e(a_U U+a_D D)P-mP.
\end{aligned}}
\label{eq:two-prey-switching}
\end{equation}
When $\mu_{UD}=\mu_{DU}=0$, this reduces to the pure sorting model. When one prey type
is removed and the carrying capacity is sent formally to infinity, the structure
reduces to the classical Lotka--Volterra predator--prey feedback.

\subsection{Spatial Extension}

The research question is spatial: we want to know whether pattern formation changes
the predator rescue threshold. Roy et al. provide the closest spatial precursor by
studying defended and undefended prey in a reaction--diffusion eco-evolutionary model
with Turing analysis \cite{roy2026}. The present model is deliberately simpler: it
keeps the defended--undefended prey structure but focuses on comparing predator
persistence thresholds between the ODE and PDE systems. We therefore allow each
population to move randomly in space. Let $U(x,t)$, $D(x,t)$, and $P(x,t)$ be densities on a habitat
$\Omega\subset\mathbb{R}^n$. The reaction--diffusion version of
\eqref{eq:two-prey-switching} is
\begin{equation}
\boxed{
\begin{aligned}
    \frac{\partial U}{\partial t}
    &= \delta_U\Delta U
    + r_U U\left(1-\frac{U+D}{K}\right)-a_U U P
    +\mu_{DU}D-\mu_{UD}U,\\
    \frac{\partial D}{\partial t}
    &= \delta_D\Delta D
    + r_D D\left(1-\frac{U+D}{K}\right)-a_D D P
    +\mu_{UD}U-\mu_{DU}D,\\
    \frac{\partial P}{\partial t}
    &= \delta_P\Delta P
    + e(a_U U+a_D D)P-mP.
\end{aligned}}
\label{eq:spatial-working-model}
\end{equation}
Here $\delta_U$, $\delta_D$, and $\delta_P$ are diffusion coefficients. We use the
letter $\delta$ for diffusion to avoid confusing diffusion parameters with the defended
prey variable $D$.
A closed habitat is represented by no-flux boundary conditions,
\begin{equation}
    \nabla U\cdot\mathbf{n}
    =
    \nabla D\cdot\mathbf{n}
    =
    \nabla P\cdot\mathbf{n}
    =0
    \qquad \text{on } \partial\Omega.
\end{equation}

The model \eqref{eq:spatial-working-model} is the extended Lotka--Volterra model for
this study. It keeps the predator--prey feedback of the classical equations, but adds
the four ingredients needed to test the proposed pattern-mediated rescue mechanism:
logistic resource limitation, defended and undefended prey, predator mortality stress,
and spatial movement.

\subsection{Predator Rescue Threshold}

The predator per capita growth rate in the non-spatial model is
\begin{equation}
    \frac{1}{P}\frac{dP}{dt}
    =
    e(a_U U+a_D D)-m.
\end{equation}
Thus predator persistence depends not only on total prey density, but also on prey
composition. For fixed total prey density, a community dominated by undefended prey
supports higher predator growth than a community dominated by defended prey.

This motivates the central comparison of the project. Let $m_c^{\mathrm{ODE}}$ be the
largest predator mortality for which the non-spatial system maintains predators above
an extinction threshold, and let $m_c^{\mathrm{PDE}}$ be the corresponding threshold in
the spatial system. The proposed pattern-mediated effect is measured by
\begin{equation}
    \Delta m_c=m_c^{\mathrm{PDE}}-m_c^{\mathrm{ODE}}.
\end{equation}
If $\Delta m_c>0$, spatial patterning promotes indirect rescue. If
$\Delta m_c<0$, spatial patterning inhibits rescue.

\subsection{Operational Definition of Persistence}

In deterministic ODE and PDE models, extinction is not always represented by an exact
finite-time crossing to zero. For numerical work, predator persistence must therefore
be defined operationally. For a final time $T$ and a small extinction threshold
$\varepsilon>0$, predators are considered persistent in the ODE model if
\begin{equation}
    P(T)>\varepsilon.
\end{equation}
In the PDE model, predators are considered persistent if the total predator biomass
satisfies
\begin{equation}
    \int_{\Omega}P(x,T)\,dx>\varepsilon.
\end{equation}
The thresholds $m_c^{\mathrm{ODE}}$ and $m_c^{\mathrm{PDE}}$ should be computed using
this criterion, and sensitivity to $T$ and $\varepsilon$ should be checked.

\section{Turing Pattern Formation}

\subsection{Why Diffusion Can Create Patterns}

Diffusion is usually thought of as a smoothing process. If one population is high in
one place and low in another, random movement tends to reduce the difference. This
makes Turing pattern formation initially surprising: under suitable conditions,
diffusion can destabilize a homogeneous state that would otherwise be stable
\cite{turing1952,murray2002}. The outcome is not uniform spreading, but the emergence of a
spatial pattern with a characteristic wavelength.

The key point is that diffusion acts on a system with several interacting variables. A
single diffusing population with stable local dynamics will not generate a Turing
pattern. But if one component locally amplifies a perturbation while another suppresses
it, and if the two components move at sufficiently different rates, diffusion can turn a
stable local balance into a spatially patterned state.

\subsection{The Two-Variable Calculation}

The cleanest derivation uses two variables. Consider
\begin{align}
    \frac{\partial u}{\partial t}
    &= f(u,v)+D_1\Delta u,\\
    \frac{\partial v}{\partial t}
    &= g(u,v)+D_2\Delta v.
\end{align}
Suppose the non-spatial system has a homogeneous equilibrium
$(u^*,v^*)$ satisfying
\begin{equation}
    f(u^*,v^*)=0,
    \qquad
    g(u^*,v^*)=0.
\end{equation}
Let
\begin{equation}
    J=
    \begin{pmatrix}
        f_u & f_v\\
        g_u & g_v
    \end{pmatrix}_{(u^*,v^*)}
\end{equation}
be the Jacobian of the reaction terms at the equilibrium. The equilibrium is stable in
the well-mixed ODE if
\begin{equation}
    \operatorname{tr}J=f_u+g_v<0,
    \qquad
    \det J=f_ug_v-f_vg_u>0.
\end{equation}

Now perturb the equilibrium by a spatial Fourier mode:
\begin{equation}
    u(x,t)=u^*+\widehat u(t)e^{i k\cdot x},
    \qquad
    v(x,t)=v^*+\widehat v(t)e^{i k\cdot x}.
\end{equation}
After linearization, this mode evolves according to the modified Jacobian
\begin{equation}
    J_k=
    \begin{pmatrix}
        f_u-D_1k^2 & f_v\\
        g_u & g_v-D_2k^2
    \end{pmatrix}.
\end{equation}
The trace becomes more negative:
\begin{equation}
    \operatorname{tr}J_k=\operatorname{tr}J-(D_1+D_2)k^2.
\end{equation}
The determinant, however, is
\begin{equation}
    \det J_k
    =
    \det J
    -
    k^2(D_1g_v+D_2f_u)
    +
    D_1D_2k^4.
\end{equation}
Thus diffusion can destabilize the equilibrium if $\det J_k<0$ for some nonzero
wavenumber $k$. Equivalently, the quadratic in $k^2$ must dip below zero. A standard
form of the Turing condition is
\begin{equation}
    D_1g_v+D_2f_u>0,
    \qquad
    (D_1g_v+D_2f_u)^2>4D_1D_2\det J.
\end{equation}
These inequalities say that the equilibrium is stable to homogeneous perturbations
but unstable to some spatially heterogeneous perturbations.

\subsection{Different Diffusion Rates}

The instability requires sufficiently different diffusion rates. If $D_1=D_2=D$, then
\begin{equation}
    J_k=J-Dk^2I.
\end{equation}
The eigenvalues are simply shifted by $-Dk^2$, so diffusion only makes the stable
equilibrium more stable. Turing instability needs differential movement. In the
standard activator--inhibitor interpretation, one component locally amplifies a
perturbation, while another component suppresses it and diffuses faster. Local
activation combined with long-range inhibition can produce spatial structure.

\subsection{Pattern Wavelength and Boundaries}

The linear calculation identifies which spatial modes grow. The fastest-growing mode
has wavenumber $k_{\max}$ and gives an expected pattern length scale
\begin{equation}
    \ell_*\approx \frac{2\pi}{k_{\max}}.
\end{equation}
This does not by itself determine the final pattern. Spots, stripes, labyrinths, or
patchy structures depend on the nonlinear dynamics after the instability grows.

On a finite domain with no-flux boundary conditions, the allowed values of $k^2$ are
not continuous. They are eigenvalues of the Laplacian on the habitat. A pattern appears
only if one or more allowed modes lie in the unstable window. This matters for
ecology: habitat size and shape can determine whether a spatial instability is visible,
and whether the resulting pattern resembles stripes, spots, or isolated patches.

\subsection{How This Enters the Rescue Problem}

The two-variable calculation is useful because it shows the mechanism clearly. Our
working model has three interacting variables, so the Turing conditions no longer
reduce to the simple trace--determinant inequalities above. In practice, the stability
test is spectral: for each admissible spatial mode, inspect the eigenvalues of the
diffusion-modified Jacobian.

For our model, define
\begin{equation}
    \mathbf{w}=(U,D,P)^T
\end{equation}
and let $\mathbf{F}(\mathbf{w})$ be the reaction part of
\eqref{eq:spatial-working-model}.
At a homogeneous equilibrium $\mathbf{w}^*$, the Jacobian of the reaction terms is
\begin{equation}
    J=\left.\frac{\partial\mathbf{F}}{\partial\mathbf{w}}\right|_{\mathbf{w}^*}.
\end{equation}
A spatial perturbation of wavenumber $k$ satisfies
\begin{equation}
    \lambda\mathbf{q}
    =
    \left[
    J-k^2\operatorname{diag}(\delta_U,\delta_D,\delta_P)
    \right]\mathbf{q},
\end{equation}
or equivalently
\begin{equation}
    \det\left(
    J-k^2\operatorname{diag}(\delta_U,\delta_D,\delta_P)-\lambda I
    \right)=0.
\end{equation}
A Turing instability occurs when all eigenvalues of $J$ have negative real parts, but
there exists some $k>0$ for which
$J-k^2\operatorname{diag}(\delta_U,\delta_D,\delta_P)$ has an eigenvalue with positive
real part.

This is the mathematical bridge to the biological question. Spatial patterning may
create local patches where predators and undefended prey overlap, raising predator
growth. It may also separate predators from edible prey, lowering predator growth.
Turing analysis is therefore not the final result. It identifies whether spatial
patterns can arise; the rescue question is answered only after comparing
$m_c^{\mathrm{PDE}}$ with $m_c^{\mathrm{ODE}}$.

\section{Planned Analysis}

\begin{enumerate}[leftmargin=2.2em]
    \item \textbf{Non-spatial baseline.} Set
    $\delta_U=\delta_D=\delta_P=0$. Find homogeneous equilibria and determine the
    predator persistence threshold $m_c^{\mathrm{ODE}}$.

    \item \textbf{Local stability.} Linearize the reaction terms around a positive
    homogeneous equilibrium and compute the Jacobian $J$.

    \item \textbf{Turing instability.} Add diffusion and scan the eigenvalues of
    \begin{equation}
        J-k^2\operatorname{diag}(\delta_U,\delta_D,\delta_P)
    \end{equation}
    over admissible spatial modes $k$. Identify parameter regions where the ODE
    equilibrium is stable but spatial modes are unstable.

    \item \textbf{Spatial threshold.} Simulate the PDE model with no-flux boundary
    conditions and compute $m_c^{\mathrm{PDE}}$.

    \item \textbf{Threshold comparison.} Evaluate
    \begin{equation}
        \Delta m_c=m_c^{\mathrm{PDE}}-m_c^{\mathrm{ODE}}.
    \end{equation}
    Use this to classify whether spatial patterning promotes rescue, inhibits rescue,
    or has a parameter-dependent effect.

    \item \textbf{Mechanistic diagnostics.} Measure total predator biomass, spatial
    overlap between predators and undefended prey, and mean edible-prey availability
    \begin{equation}
        \langle a_UU+a_DD\rangle.
    \end{equation}
    These diagnostics should explain why $\Delta m_c$ is positive or negative.

    \item \textbf{Robustness checks.} Vary initial prey composition, extinction
    threshold, final time $T$, domain size, diffusion ratios, conversion rates, and at
    least one saturating functional response such as Holling type II predation.
\end{enumerate}

\section{Expected Contributions and Limitations}

The expected contribution is a focused threshold comparison. Rather than claiming that
pattern-mediated indirect evolutionary rescue is already established, the study tests
whether a spatially patterned predator--prey system can support predators at mortality
levels that would cause loss in the corresponding non-spatial model.
The main scientific outcome is not the existence of a spatial pattern itself, but
whether that pattern changes the predator persistence threshold relative to the
well-mixed system.

There are three possible outcomes. If $\Delta m_c>0$, spatial patterning expands the
predator persistence window and acts as a rescue-promoting mechanism. If
$\Delta m_c<0$, patterning separates predators from edible prey or otherwise lowers
effective predator growth, so spatial structure inhibits rescue. If $\Delta m_c$
changes sign across parameter regimes, the main result is conditional: pattern
formation can either help or harm predator persistence depending on diffusion rates,
prey defence costs, and spatial overlap.

The main limitation is that the conversion terms between $U$ and $D$ are
phenomenological. They make coexistence and Turing analysis mathematically well posed,
but they do not by themselves identify whether the biological source is mutation,
plastic switching, inducible defence, or persistent standing variation. A later model
could separate these mechanisms once the minimal threshold question is understood.
Similarly, the baseline model uses mass-action predation; robustness checks with a
saturating functional response are needed before treating the threshold shift as a
general ecological prediction.

\begin{thebibliography}{99}

\bibitem{turing1952}
Turing, A. M. (1952).
The chemical basis of morphogenesis.
\emph{Philosophical Transactions of the Royal Society of London. Series B, Biological
Sciences}, 237(641), 37--72.
\url{https://doi.org/10.1098/rstb.1952.0012}.

\bibitem{gomulkiewicz1995}
Gomulkiewicz, R., and Holt, R. D. (1995).
When does evolution by natural selection prevent extinction?
\emph{Evolution}, 49(1), 201--207.

\bibitem{yamamichi2015}
Yamamichi, M., and Miner, B. E. (2015).
Indirect evolutionary rescue: prey adapts, predator avoids extinction.
\emph{Evolutionary Applications}, 8(8), 787--795.

\bibitem{hermann2022}
Hermann, R. J., and Becks, L. (2022).
Change in prey genotype frequency rescues predator from extinction.
\emph{Royal Society Open Science}, 9, 220211.
\url{https://doi.org/10.1098/rsos.220211}.

\bibitem{kasada2014}
Kasada, M., Yamamichi, M., and Yoshida, T. (2014).
Form of an evolutionary tradeoff affects eco-evolutionary dynamics in a predator--prey
system.
\emph{Proceedings of the National Academy of Sciences}, 111(45), 16035--16040.
\url{https://doi.org/10.1073/pnas.1406357111}.

\bibitem{hermann2024}
Hermann, R. J., Pantel, J. H., R{\'e}veillon, T., and Becks, L. (2024).
Range of trait variation in prey determines evolutionary contributions to predator
growth rates.
\emph{Journal of Evolutionary Biology}, 37, 693--703.
\url{https://doi.org/10.1093/jeb/voae062}.

\bibitem{roy2026}
Roy, S., Byrne, H. M., Banerjee, J., Sasmal, S. K., Dutta, S., and Ghosh, D. (2026).
Spatio-temporal eco-evolutionary dynamics of prey-predator systems with defended and
undefended prey.
\emph{Communications Physics}, 9, 4.
\url{https://doi.org/10.1038/s42005-025-02434-1}.

\bibitem{murray2002}
Murray, J. D. (2002).
\emph{Mathematical Biology I: An Introduction}.
Third edition. Interdisciplinary Applied Mathematics, Vol. 17. Springer.
\url{https://doi.org/10.1007/b98868}.

\bibitem{kot2001}
Kot, M. (2001).
\emph{Elements of Mathematical Ecology}.
Cambridge University Press.
\url{https://doi.org/10.1017/CBO9780511608520}.

\bibitem{edelstein2005}
Edelstein-Keshet, L. (2005).
\emph{Mathematical Models in Biology}.
Classics in Applied Mathematics. SIAM.
\url{https://doi.org/10.1137/1.9780898719147}.

\end{thebibliography}

\end{document}
